\section{Empirical analysis of the cycle}

\begin{frame}[t]
  \frametitle{Organic composition of capital and rate of profit (1)\footnotetext{Source: FRED, NBER, BEA}}
  \begin{center}
    \input{Plots/occ-profit.tex} % Органическое строение капитала и норма прибыли
  \end{center}
\end{frame}

\begin{frame}
  \frametitle{Organic composition of capital and rate of profit (2)}
  \begin{block}{Observation}
    \begin{itemize}
      \item \textbf{Organic composition of capital} grows at the end of the growth phase due to the commissioning of production facilities created as a result of previous investments; during the crisis its growth is additionally supported by a sharp decline in employment and wages.
      \item \textbf{Rate of profit} moves in the opposite direction, begins to decline from about the middle of the industrial cycle under the pressure of constant capital accumulation outpacing the growth of surplus value, during the crisis it additionally decreases due to the depreciation of capital, after which it temporarily recovers.
    \end{itemize}
  \end{block}

  \begin{alertblock}{Finding}
    Empirically obtained results on the basis of up-to-date statistical data confirm the fundamental principles of marxist political economy regarding the long-term trend of increasing the organic composition of capital and decreasing the rate of profit.
  \end{alertblock}
\end{frame}

\begin{frame}[t]
  \frametitle{Fixed capital average age (1)\footnotetext{Source: FRED, NBER, BEA}}
  \begin{center}
    \input{Plots/average-age.tex} % Средний возраст основного капитала
  \end{center}
\end{frame}

\begin{frame}
  \frametitle{Fixed capital average age (2)}
  \begin{block}{Observation}
    \textbf{Average age of fixed capital} decreases at the beginning of the growth phase during technological modernization and expansion of production capacities, and increases during the crisis due to the suspension of large investments.
  \end{block}

  \begin{alertblock}{Finding}
    The dynamics of the average age of fixed capital confirms the propositions about the restoration of production during the recovery phase due to existing production capacities and the massive renewal of fixed assets at the initial stage of the growth phase.
  \end{alertblock}
\end{frame}

\begin{frame}[t]
  \frametitle{Private fixed investment (1)\footnotetext{Source: FRED, NBER}}
  \begin{center}
    \input{Plots/private-fixed-investment.tex} % Частные инвестиции в основной капитал
  \end{center}
\end{frame}

\begin{frame}
  \frametitle{Private fixed investment (2)}
  \begin{block}{Observation}
    \textbf{Private fixed investment} generally follows the dynamics of the industrial production index, but at the end of the growth phase begin to decline earlier, and after the recovery phase they recover a little later and more slowly.
  \end{block}

  \begin{alertblock}{Finding}
    The dynamics of private fixed investment confirms the proposition about the recovery of production due to existing production capacities; a reduction in investment at the end of the growth phase signals the approach of a cyclical crisis of general overproduction.
  \end{alertblock}
\end{frame}

\begin{frame}[t]
  \frametitle{Manufacturers' new orders (1)\footnotetext{Source: FRED, NBER}}
  \begin{center}
    \input{Plots/orders.tex} % Портфель заказов
  \end{center}
\end{frame}

\begin{frame}
  \frametitle{Manufacturers' new orders (2)}
  \begin{block}{Observation}
    \textbf{Manufacturers' new orders} generally follow the dynamics of private fixed investment, shrink earlier at the end of the growth phase, and recover with some delay after the recovery phase.
  \end{block}

  \begin{alertblock}{Finding}
    The slump of the manufacturers' new orders at the end of the growth phase signals the approach of a cyclical crisis of general overproduction.
  \end{alertblock}
\end{frame}

\begin{frame}
  \frametitle{Capacity utilization (1)\footnotetext{Source: FRED, NBER}}
  \begin{center}
    \input{Plots/capacity.tex} % Загрузка производственных мощностей
  \end{center}
\end{frame}

\begin{frame}
  \frametitle{Capacity utilization (2)}
  \begin{block}{Observation}
    \textbf{Capacity utilization} exactly repeats the dynamics of the industrial production index and changes synchronously with it, decreases during the crisis, increases during the recovery and growth.
  \end{block}

  \begin{alertblock}{Finding}
    Capacity utilization can be considered solely as a current indicator of the industrial cycle in order to concretize the assessment of the depth of the decline in production volumes.
  \end{alertblock}
\end{frame}

\begin{frame}
  \frametitle{Inventories (1)\footnotetext{Source: FRED, NBER}}
  \begin{center}
    \input{Plots/inventories.tex} % Товарно-материальные запасы
  \end{center}
\end{frame}

\begin{frame}
  \frametitle{Inventories (2)}
  \begin{block}{Observation}
    \textbf{Inventories} decrease during the crisis due to a slump in production and the desire of enterprises to free up associated working capital, then the rate of their accumulation increases until about the middle of the growth phase under the influence of expanding output and expectations of further growth in demand, but slows down towards the end due to the sales difficulties.
  \end{block}

  \begin{alertblock}{Finding}
    The dynamics of inventories clearly illustrates the aggravation and temporary resolution of the contradiction between the conditions of production and the conditions of realization of surplus value, as well as the process of depreciation of the commodity form of capital during the crisis of overproduction.
  \end{alertblock}
\end{frame}

\begin{frame}
  \frametitle{Producer and consumer prices (1)\footnotetext{Source: FRED, NBER}}
  \begin{center}
    \input{Plots/prices.tex} % Цены производителей и потребителей
  \end{center}
\end{frame}

\begin{frame}
  \frametitle{Producer and consumer prices (2)}
  \begin{block}{Observation}
    \begin{itemize}
      \item The growth rate of the \textbf{producer price index} increases at the beginning of the crisis phase under the influence of previously accumulated cost increases and inertia of contract prices, but as the recession deepens it sharply decreases due to reduced demand and difficulties in sales, and then recovers at the beginning of the growth phase.
      \item The growth rate of the \textbf{consumer price index} generally follows the dynamics of the producer price index, but is less volatile and rarely falls below zero under the influence of inflation.
    \end{itemize}
  \end{block}

  \begin{alertblock}{Finding}
    The dynamics of the growth rates of the producer and consumer price indices clearly illustrate the process of depreciation of the commodity form of capital during the crisis of overproduction.
  \end{alertblock}
\end{frame}

\begin{frame}
  \frametitle{Unemployment (1)\footnotetext{Source: FRED, NBER}}
  \begin{center}
    \input{Plots/unemployment.tex} % Безработица
  \end{center}
\end{frame}

\begin{frame}
  \frametitle{Unemployment (2)}
  \begin{block}{Observation}
    \textbf{Unemployment} increases sharply at the beginning of the crisis phase under the influence of a slump in production and mass release of labor, starting from the recovery phase employment increases as workers are involved in the production process and wages increase, and by the end of the growth phase the unemployment rate reaches a minimum.
  \end{block}
  
  \begin{alertblock}{Finding}
    The dynamics of the unemployment rate clearly illustrates the steady trend of relative overpopulation within the capitalist economy and the replenishment of the reserve army of labor in the context of the crisis of overproduction.
  \end{alertblock}
\end{frame}

\begin{frame}
  \frametitle{Money supply and velocity of circulation (1)\footnotetext{Source: FRED, NBER}}
  \begin{center}
    \input{Plots/m2-velocity.tex} % Денежная масса и скорость обращения денег
  \end{center}
\end{frame}

\begin{frame}
  \frametitle{Money supply and velocity of circulation (2)}
  \begin{block}{Observation}
    \begin{itemize}
      \item \textbf{Money supply} during a crisis can either shrink under the influence of a squeeze on bank lending and an raise in non-payments or increase in the case of a stimulating monetary policy.
      \item \textbf{Velocity of circulation} varies inversely with money supply, but always decreases during a crisis under the influence of reduced investment and disruption of payment chains.
    \end{itemize}
  \end{block}

  \begin{alertblock}{Finding}
    The dynamics of the money supply and velocity of circulation clearly illustrate the process of transformation of the credit system into a monetary one and the growing demand for money as a means of payment in the context of the crisis of overproduction.
  \end{alertblock}
\end{frame}

\begin{frame}
  \frametitle{Stock market (1)\footnotetext{Source: FRED, NBER, Yahoo Finance}}
  \begin{center}
    \input{Plots/stock-market.tex} % Фондовый рынок
  \end{center}
\end{frame}

\begin{frame}
  \frametitle{Stock market (2)}
  \begin{block}{Observation}
    \textbf{Stock market indices} generally repeat the dynamics of the industrial production index, characterized by an earlier decline at the end of the growth phase since the pricing of financial assets is based on the level of expected income.
  \end{block}

  \begin{alertblock}{Finding}
    The dynamics of stock market indices reflect the increased role of fictitious capital in the context of expanded capitalist reproduction under the domination of financial capital as well as the process of depreciation of financial capital in the context of the crisis of overproduction.
  \end{alertblock}
\end{frame}

\begin{frame}
  \frametitle{Net exports (1)\footnotetext{Source: FRED, NBER}}
  \begin{center}
    \input{Plots/net-export.tex} % Чистый экспорт
  \end{center}
\end{frame}

\begin{frame}
  \frametitle{Net exports (2)}
  \begin{block}{Observation}
    \textbf{Net exports} in the US economy have a long-term downward trend, but during the crisis phase temporarily recover under the influence of a reduction in national income and import volumes.
  \end{block}

  \begin{alertblock}{Finding}
    The dynamics of net exports clearly illustrates the compression of export potential, domestic market and economic turnover relative to other countries during the crisis of overproduction.
  \end{alertblock}
\end{frame}

\begin{frame}
  \frametitle{Foreign direct investment (1)\footnotetext{Source: FRED, NBER, UNCTAD}}
  \begin{center}
    \input{Plots/fdi-flow.tex} % Прямые иностранные инвестиции
  \end{center}
\end{frame}

\begin{frame}
  \frametitle{Foreign direct investment (2)}
  \begin{block}{Observation}
    Outward and inward \textbf{foreign direct investment} generally follow the dynamics of private fixed investment: decrease earlier at the end of the growth phase and recover a little later after the recovery phase.
  \end{block}

  \begin{alertblock}{Finding}
    The dynamics of outward and inward foreign direct investment vividly illustrates international capital flows in relation to the industrial cycle at the national level; a decrease in foreign direct investment at the end of the growth phase can be considered as a leading indicator of the crisis of overproduction.
  \end{alertblock}
\end{frame}

\begin{frame}
  \frametitle{Government expenditures and interest rate (1)\footnotetext{Source: FRED, NBER}}
  \begin{center}
    \input{Plots/gov-consumption-key-rate.tex} % Государственные расходы и ключевая ставка
  \end{center}
\end{frame}

\begin{frame}
  \frametitle{Government expenditures and interest rate (2)}
  \begin{block}{Observation}
    \begin{itemize}
      \item \textbf{Government expenditures} in the US economy have a long-term upward trend, but during the crisis phase their growth rates increase under the influence of stimulating government policies.
      \item \textbf{Interest rate} increases towards the end of the growth phase to limit inflationary pressures and decreases during the crisis in order to support bank liquidity and credit growth which corresponds to the transition from a restraining to a stimulating monetary policy.
    \end{itemize}
  \end{block}

  \begin{alertblock}{Finding}
    The dynamics of government expenditures and the central bank's interest rate clearly illustrate the nature of government regulatory policy in the context of the phases of the industrial cycle.
  \end{alertblock}
\end{frame}